\documentclass[11pt,letterpaper]{article}
\usepackage[utf8]{inputenc}
\usepackage[T1]{fontenc}
\usepackage[spanish]{babel}
\usepackage{amsmath,amssymb}
\usepackage[a4paper,margin=2.4cm]{geometry}
\usepackage{enumitem}
\usepackage{hyperref}
\setlength{\parindent}{0pt}
\setlength{\parskip}{6pt}

\title{Transcripcion parcial de fotos}
\author{Apuntes de Modelos Matematicos I}
\date{}

\begin{document}
\maketitle

\begin{quote}
Nota: esta es una transcripcion parcial y acumulativa. Las fotos de baja resolucion
se leyeron con varias pasadas; donde el texto era dudoso se dejaron marcas de
incertidumbre o se priorizo la formula legible.
\end{quote}

\section*{Foto 1}
\texttt{Fotos/photo\_2026-06-04\_00-09-12.jpg}

\textbf{Contexto probable:} problema de juguete, ecuacion autonoma escalar, campo de pendientes y estabilidad del equilibrio.

\textbf{Transcripcion tentativa:}
\begin{itemize}[leftmargin=*]
    \item ``Problema de juguete''.
    \item [Texto superior izquierdo muy borroso; solo se distinguen fragmentos como ``Decimos que...''].
    \item \[
    \frac{du}{dt}=f(u), \qquad \frac{d\mu}{dt}=f(\mu)
    \]
    \item Caso 1:
    \[
    \frac{du}{dt}=ru, \qquad r\in\mathbb{R}
    \]
    \item Se dibujan las graficas de \(f(u)=ru\) para \(r>0\), \(r<0\) y \(r=0\), junto con el campo de pendientes y la linea de fase.
    \item Idea que se repite en el dibujo: si \(r>0\), el equilibrio \(u=0\) es inestable; si \(r<0\), es estable; si \(r=0\), el criterio lineal no decide.
\end{itemize}

\section*{Foto 2}
\texttt{Fotos/photo\_2026-06-04\_00-09-39.jpg}

\textbf{Contexto probable:} linealizacion de sistemas en dos variables, puntos de equilibrio y jacobiano.

\textbf{Transcripcion tentativa:}
\begin{itemize}[leftmargin=*]
    \item \[
    x' = F(x), \qquad x_0 \text{ pto. de equilibrio}, \qquad F(x_0)=0.
    \]
    \item ``Calculamos el jacobiano de''
    \[
    F=(F_1(x,y),F_2(x,y)).
    \]
    \item
    \[
    J(x)=
    \begin{pmatrix}
    \dfrac{\partial F_1}{\partial x}(x,y) & \dfrac{\partial F_1}{\partial y}(x,y)\\[6pt]
    \dfrac{\partial F_2}{\partial x}(x,y) & \dfrac{\partial F_2}{\partial y}(x,y)
    \end{pmatrix}
    \]
    \item Ejemplo:
    \[
    x' = x^2-y^2-1, \qquad y' = 2y
    \]
    \item
    \[
    DF(x)=J(x)=
    \begin{pmatrix}
    2x & -2y\\
    0 & 2
    \end{pmatrix}
    \]
    \item Puntos de equilibrio:
    \[
    x^2-y^2-1=0, \qquad 2y=0
    \]
    \item Como \(y=0\), entonces \(x^2-1=0\), y por tanto \(x=\pm 1\).
    \item Se concluye:
    \[
    X_0=(1,0), \qquad X_1=(-1,0).
    \]
\end{itemize}

\section*{Foto 3}
\texttt{Fotos/photo\_2026-06-04\_00-09-25.jpg}

\textbf{Contexto probable:} reescalamiento del modelo logistico para llevarlo a la forma normalizada.

\textbf{Transcripcion tentativa:}
\begin{itemize}[leftmargin=*]
    \item \[
    \frac{dU}{dt}=rU\left(1-\frac{U}{n}\right)=rU-\frac{r}{n}U^2
    \]
    \item ``Dos puntos de equilibrio'':
    \[
    U_0=0, \qquad U_1=n
    \]
    \item ``Hagamos: cambio de variable espacial y temporal''.
    \item Definicion de la variable adimensional:
    \[
    M=\phi(U)=\frac{U}{n}, \qquad \phi^{-1}(M)=nM
    \]
    \item Reescritura con cambio temporal:
    \[
    \tau=rt, \qquad t=\frac{\tau}{r}, \qquad \frac{dt}{d\tau}=\frac{1}{r}
    \]
    \item Forma normalizada:
    \[
    \frac{dM}{d\tau}=M(1-M)
    \]
    \item La pregunta escrita al centro parece ser algo como: ``¿Por que se elige el cambio temporal?''.
\end{itemize}

\section*{Foto 4}
\texttt{Fotos/photo\_2026-06-04\_00-09-16.jpg}

\textbf{Contexto probable:} modelo logistico, crecimiento intrinseco y cosecha constante.

\textbf{Transcripcion tentativa:}
\begin{itemize}[leftmargin=*]
    \item Titulo: ``Modelo logistico''.
    \item Forma clasica:
    \[
    \frac{dN}{dt}=rN\left(1-\frac{N}{u}\right)
    \]
    \item Forma cuadratica equivalente:
    \[
    \frac{dN}{dt}=\alpha N+\beta N^2
    \]
    \item ``Tasa de crecimiento intrinseca''.
    \item Reescritura:
    \[
    \frac{dN}{dt}=\alpha N\left(1-\frac{\beta}{\alpha}N\right)
    =\alpha N\left(1-\frac{N}{K}\right)
    \]
    \item Hay esquemas de ``entrada/salida'' y una interpretacion de \( \beta N^2 \) como termino de salida.
    \item En la parte derecha aparece el modelo con cosecha:
    \[
    \frac{dN}{dt}=rN\left(1-\frac{N}{K}\right)-H
    \]
    \item Puntos criticos:
    \[
    rN\left(1-\frac{N}{K}\right)-H=0
    \]
    \[
    -\frac{r}{K}N^2+rN-H=0
    \]
    \item Solucion cuadratica:
    \[
    N=\frac{K}{2}\left(1\pm \sqrt{1-\frac{4H}{rK}}\right)
    \]
    \item Casos anotados:
    \begin{itemize}
        \item un unico equilibrio si \(H=\frac{rK}{4}\),
        \item no hay equilibrios reales si \(H>\frac{rK}{4}\),
        \item dos equilibrios reales si \(H<\frac{rK}{4}\).
    \end{itemize}
\end{itemize}

\section*{Foto 5}
\texttt{Fotos/photo\_2026-06-04\_00-09-18.jpg}

\textbf{Contexto probable:} normalizacion del modelo logistico y observaciones sobre el cambio de variables.

\textbf{Transcripcion tentativa:}
\begin{itemize}[leftmargin=*]
    \item Se repite la ecuacion logistico-normalizada:
    \[
    \frac{dN}{dt}=rN\left(1-\frac{N}{n}\right)
    \]
    \item ``Dos puntos de equilibrio'':
    \[
    N_0=0, \qquad N_1=n
    \]
    \item ``Hagamos: cambio de variable espacial y temporal''.
    \item Cambio de variable:
    \[
    M=\phi(N)=\frac{N}{n}, \qquad \phi^{-1}(M)=nM
    \]
    \item Relacion entre derivadas:
    \[
    \frac{dN}{dt}=n\frac{dM}{dt}
    \]
    \item Con el cambio temporal:
    \[
    \tau=rt, \qquad \frac{dt}{d\tau}=\frac{1}{r}
    \]
    \item Forma adimensional:
    \[
    \frac{dM}{d\tau}=M(1-M)
    \]
    \item La nota al centro deja ver la idea de que \(M\) y \(N\) representan la misma solucion en escalas distintas.
\end{itemize}

\section*{Foto 6}
\texttt{Fotos/photo\_2026-06-04\_00-09-20.jpg}

\textbf{Contexto probable:} forma logistica desde el modelo cuadratico, y analisis de la ecuacion con cosecha.

\textbf{Transcripcion tentativa:}
\begin{itemize}[leftmargin=*]
    \item \[
    \frac{dN}{dt}=rN\left(1-\frac{N}{K}\right)
    \]
    \item \[
    \frac{dN}{dt}=\alpha N+\beta N^2
    \]
    \item ``Tasa de crecimiento intrinseca''.
    \item Reescrituras:
    \[
    =\alpha N\left(1-\frac{\beta}{\alpha}N\right)
    =\alpha N\left(1-\frac{N}{K}\right)
    \]
    \item En el dibujo se interpreta un proceso de ``entrada'' y ``salida'' con un termino no lineal proporcional a \(N^2\).
    \item Hacia la derecha aparece el modelo con cosecha:
    \[
    \frac{dN}{dt}=rN\left(1-\frac{N}{K}\right)-H
    \]
    \item Se dibuja la parabola \(f(N)=rN(1-N/K)-H\) y varias lineas de fase para los casos \(a=0\), \(a>0\), \(a<0\) de una bifurcacion tipo silla-nodo.
\end{itemize}

\section*{Foto 7}
\texttt{Fotos/photo\_2026-06-04\_00-09-23.jpg}

\textbf{Contexto probable:} analisis biologico del modelo logistico con cosecha constante.

\textbf{Transcripcion tentativa:}
\begin{itemize}[leftmargin=*]
    \item Se trabaja con
    \[
    f(N)=rN\left(1-\frac{N}{K}\right)-H
    =-\frac{r}{K}N^2+rN-H
    \]
    \item Condicion de dos equilibrios:
    \[
    H<\frac{rK}{4}
    \]
    \item La grafica muestra una parabola hacia abajo y una linea de fase con un equilibrio inestable y otro estable.
    \item ``Analisis biologico''.
    \item La lectura global es:
    \begin{itemize}
        \item si la poblacion inicial esta por encima del umbral critico, la solucion tiende al equilibrio estable superior;
        \item si la poblacion inicial esta por debajo del umbral, la poblacion puede extinguirse en tiempo finito;
        \item el equilibrio inferior actua como separatriz o punto umbral.
    \end{itemize}
    \item En la esquina derecha se alcanza a leer ``2\(^\circ\) caso'' y ``3\(^\circ\) caso'', lo que sugiere que se estaban enumerando escenarios segun la condicion inicial.
\end{itemize}

\section*{Foto 8}
\texttt{Fotos/photo\_2026-06-04\_00-09-28.jpg}

\textbf{Contexto probable:} analisis dimensional: unidades, dimensiones y ejemplo con presion.

\textbf{Transcripcion tentativa:}
\begin{itemize}[leftmargin=*]
    \item ``Observamos un modelo cuyos resultados sean adimensionales del sistema de unidades''.
    \item ``Para ello hacemos:''
    \begin{enumerate}[leftmargin=*]
        \item Unidades basicas del SI: metro (m), kilogramo (kg), segundo (s), ampere (A), mol (mol), candela (cd).
        \item Unidad derivada: velocidad \(m/s\), fuerza \(N\).
        \item Dimensiones: representan la estructura de una magnitud o fenomeno. Longitud \([L]\), tiempo \([T]\), masa \([M]\).
    \end{enumerate}
    \item Tabla de unidades basicas y derivadas.
    \item Ejemplo:
    \[
    \text{Presion}=\frac{\text{Fuerza}}{\text{Area}}
    \]
    \[
    F=ma \sim \frac{kg\cdot m}{s^2}, \qquad A=m^2
    \]
    \item Dimensiones:
    \[
    [F]=[M][L][T]^{-2}
    \]
    \[
    [P]=[M][L]^{-1}[T]^{-2}
    \]
    \item Nota final visible: ``Tarea: Como se define en N o Joule?''.
\end{itemize}

\section*{Foto 9}
\texttt{Fotos/photo\_2026-06-04\_00-13-25.jpg}

\textbf{Contexto probable:} graficas de soluciones y linea de fase del modelo logístico adimensional con cosecha.

\textbf{Transcripcion tentativa:}
\begin{itemize}[leftmargin=*]
    \item Se observan curvas de solucion con forma sigmoidea y varias trayectorias de fase.
    \item Aparecen etiquetas en la pizarra:
    \[
    N=K, \qquad N=\frac{K}{2}, \qquad N=0, \qquad N_0=N(0)
    \]
    \item En la parte inferior se distingue la funcion de fase:
    \[
    f(M)=M(1-M)
    \]
    \item Tambien se ve una linea de fase con los puntos \(M=0\) y \(M=1\).
    \item La lectura general es la del modelo logistico reescalado, con soluciones que se acercan a los equilibrios visibles en la grafica.
\end{itemize}

\section*{Foto 10}
\texttt{Fotos/photo\_2026-06-04\_00-13-27.jpg}

\textbf{Contexto probable:} deduccion del modelo adimensional con cosecha y definicion del parametro \(\lambda\).

\textbf{Transcripcion tentativa:}
\begin{itemize}[leftmargin=*]
    \item Cambio temporal:
    \[
    t=\frac{\tau}{r}, \qquad \tau=rt
    \]
    \item Relacion de derivadas:
    \[
    \frac{dM}{d\tau}=\frac{dM}{dt}\frac{dt}{d\tau}
    =\frac{1}{r}\frac{dM}{dt}
    \]
    \item Sustitucion del modelo:
    \[
    \frac{dM}{d\tau}
    =\frac{1}{r}\left[rM(1-M)-\frac{H}{K}\right]
    \]
    \item Resultado final:
    \[
    \frac{dM}{d\tau}=M(1-M)-\frac{H}{rK}
    \]
    \item Se introduce
    \[
    \lambda=\frac{H}{rK}
    \]
    y por tanto
    \[
    \frac{dM}{d\tau}=M(1-M)-\lambda.
    \]
    \item Nota al margen: ``\textit{Voila!}''
    \item Tarea escrita en verde:
    \begin{quote}
        ``Realizar el analisis completo semejante al realizado sin adimensionalizacion.''
    \end{quote}
\end{itemize}

\section*{Foto 11}
\texttt{Fotos/photo\_2026-06-04\_00-09-26.jpg}

\textbf{Contexto probable:} adimensionalizacion del modelo logistico con cosecha y motivacion de por que conviene reescalar.

\textbf{Transcripcion tentativa:}
\begin{itemize}[leftmargin=*]
    \item ``Adimensionalizacion: Herramientas del proceso de modelacion''.
    \item ``Objetivo: En los modelos vistos hasta el momento aparecen parametros''.
    \item En el modelo logistico con cosecha:
    \[
    \frac{dN}{dt}=rN\left(1-\frac{N}{K}\right)-H
    \]
    \item La dinamica depende de \(H\) y aparece el umbral
    \[
    H<\frac{rK}{4}.
    \]
    \item Se plantea el cambio de variable
    \[
    u=\frac{N}{K}
    \]
    y el cambio temporal
    \[
    \tau=rt.
    \]
    \item El modelo adimensional queda, en esencia,
    \[
    \frac{du}{d\tau}=u(1-u)-\lambda,
    \qquad \lambda=\frac{H}{rK}.
    \]
    \item Idea central escrita en el pizarron: ``El nuevo modelo no depende de parametros de escala''.
    \item Abajo se aclara que la pregunta es por que se hace un reescalamiento.
\end{itemize}

\section*{Foto 12}
\texttt{Fotos/photo\_2026-06-04\_00-09-34.jpg}

\textbf{Contexto probable:} dinamica de dos poblaciones y modelo de Lotka-Volterra con respuesta funcional.

\textbf{Transcripcion tentativa:}
\begin{itemize}[leftmargin=*]
    \item ``Los niveles cero'' y lineas de fase para \(f(y)=ry\).
    \item Si \(r>0\), entonces \(N_0=0\) es inestable.
    \item Si \(r<0\), entonces \(N_0=0\) es estable.
    \item Se introduce ``Dinamica de dos poblaciones''.
    \item Modelo base de crecimiento logistico:
    \[
    \frac{dN}{dt}=rN\left(1-\frac{N}{K}\right)
    \]
    \item Con competencia / efecto Allee:
    \[
    \frac{dN}{dt}=rN\left(1-\frac{N}{K}\right)\left(\frac{N}{A}-1\right)-H
    \]
    \[
    \frac{dN}{dt}=rN\left(1-\frac{N}{K}\right)\left(\frac{N}{A}-1\right)-qEN
    \]
    \item Se menciona el modelo presa-depredador:
    \[
    \frac{dP}{dt}=\alpha P-\beta PD,
    \qquad
    \frac{dD}{dt}=-\gamma D+\delta PD.
    \]
    \item Tambien aparece la version con crecimiento logistico de la presa:
    \[
    \frac{dP}{dt}=rP\left(1-\frac{P}{K}\right)-\beta PD,
    \qquad
    \frac{dD}{dt}=-\gamma D+\delta PD.
    \]
    \item El dibujo superior asocia \(f(y)=ry\) con estabilidad del origen dependiendo del signo de \(r\).
\end{itemize}

\section*{Foto 13}
\texttt{Fotos/photo\_2026-06-04\_00-09-36.jpg}

\textbf{Contexto probable:} sistemas autonomos en el plano y sistemas de tipo Kolmogorov.

\textbf{Transcripcion tentativa:}
\begin{itemize}[leftmargin=*]
    \item ``Sistemas de ecuaciones diferenciales autonomos''.
    \item Forma general en el plano:
    \[
    \frac{dx}{dt}=P(x,y), \qquad \frac{dy}{dt}=Q(x,y).
    \]
    \item Se enfatiza que esto asociado a un campo vectorial en el plano
    \[
    F:\mathbb{R}^2 \to \mathbb{R}^2, \qquad F(x,y)=(P(x,y),Q(x,y)).
    \]
    \item Forma de un sistema de tipo Kolmogorov:
    \[
    \frac{dx}{dt}=x\,f_1(x,y), \qquad \frac{dy}{dt}=y\,f_2(x,y).
    \]
    \item En la pizarra se remarca que reflejan una caracteristica biologica importante.
    \item Se dibujan retratos de fase con equilibrio tipo silla, foco/atractor y repulsor, para ilustrar la variedad de comportamientos.
\end{itemize}

\section*{Foto 14}
\texttt{Fotos/photo\_2026-06-04\_00-09-38.jpg}

\textbf{Contexto probable:} dinamica local en sistemas no lineales y dos posibilidades: punto regular o equilibrio.

\textbf{Transcripcion tentativa:}
\begin{itemize}[leftmargin=*]
    \item ``El mundo en un mundo no lineal''.
    \item \[
    x' = F(x) \quad \text{en } \mathbb{R}^n,
    \qquad F:\mathbb{R}^n \to \mathbb{R}^n,\ \text{de clase } C^1.
    \]
    \item ``Problema general: Entender la dinamica global de las soluciones de \((*)\)''.
    \item ``Dinamica local: Dado \(x_0\in\mathbb{R}^n\), como es el comportamiento de las soluciones cerca de \(x_0\)?''.
    \item Posibilidades:
    \begin{enumerate}[leftmargin=*]
        \item \(x_0\) es un punto regular, \(F(x_0)\neq 0\).
        \item \(x_0\) es un punto de equilibrio, \(F(x_0)=0\).
    \end{enumerate}
    \item Se comienza el caso \(x_0\) punto regular.
\end{itemize}

\section*{Foto 15}
\texttt{Fotos/photo\_2026-06-04\_00-13-30.jpg}

\textbf{Contexto probable:} lineas de fase de una ecuacion autonoma escalar y estabilidad por linealizacion.

\textbf{Transcripcion tentativa:}
\begin{itemize}[leftmargin=*]
    \item \[
    \frac{dy}{dt}=f(y), \qquad f:\mathbb{R}\to\mathbb{R}
    \]
    \item Se dibuja la linea de fase para \(f(y_0)=0\).
    \item Para \(f(y)=ry\):
    \begin{itemize}
        \item si \(r>0\), entonces \(N_0=0\) es inestable;
        \item si \(r<0\), entonces \(N_0=0\) es estable.
    \end{itemize}
    \item Por linealizacion:
    \[
    f'(N_0)=r.
    \]
    \item Si \(r>0\), entonces \(f'(N_0)>0\) y el equilibrio es inestable.
    \item Si \(r<0\), entonces \(f'(N_0)<0\) y el equilibrio es estable.
\end{itemize}

\section*{Foto 16}
\texttt{Fotos/photo\_2026-06-04\_00-13-31.jpg}

\textbf{Contexto probable:} modelo presa-depredador y preparacion del modelo Lotka-Volterra.

\textbf{Transcripcion tentativa:}
\begin{itemize}[leftmargin=*]
    \item Titulo: ``Dinamica de dos poblaciones''.
    \item Se parte del crecimiento logistico:
    \[
    \frac{dN}{dt}=rN \ \longrightarrow\ \frac{dN}{dt}=rN\left(1-\frac{N}{K}\right).
    \]
    \item Luego se introducen modificaciones de competencia / Allee:
    \[
    \frac{dN}{dt}=rN\left(1-\frac{N}{K}\right)\left(\frac{N}{A}-1\right)-H
    \]
    \[
    \frac{dN}{dt}=rN\left(1-\frac{N}{K}\right)\left(\frac{N}{A}-1\right)-qEN.
    \]
    \item ``Modelo de Lotka-Volterra (1925, 1926)''.
    \item Variables:
    \[
    P(t) \text{ poblacion de presa al tiempo } t,
    \qquad
    D(t) \text{ poblacion de depredador}.
    \]
    \item Ecuaciones clasicas:
    \[
    \frac{dP}{dt}=\alpha P-\beta PD,
    \qquad
    \frac{dD}{dt}=-\gamma D+\delta PD.
    \]
    \item Se anota que existen distintas respuestas funcionales.
\end{itemize}

\section*{Foto 17}
\texttt{Fotos/photo\_2026-06-04\_00-13-33.jpg}

\textbf{Contexto probable:} formulaciones generales del sistema presa-depredador con respuesta funcional.

\textbf{Transcripcion tentativa:}
\begin{itemize}[leftmargin=*]
    \item El modelo logistico de la presa se modifica como
    \[
    \frac{dP}{dt}=rP\left(1-\frac{P}{K}\right)-\beta PD.
    \]
    \item La ecuacion del depredador queda
    \[
    \frac{dD}{dt}=-\gamma D+\delta PD.
    \]
    \item Se sugiere una respuesta funcional mas general \(f(P)\), de modo que
    \[
    \frac{dP}{dt}=rP\left(1-\frac{P}{K}\right)-f(P)D,
    \qquad
    \frac{dD}{dt}=-\gamma D+\delta f(P)D.
    \]
    \item En el pizarron se comparan distintos tipos de respuesta funcional con dibujos cualitativos.
\end{itemize}

\section*{Foto 18}
\texttt{Fotos/photo\_2026-06-04\_00-13-34.jpg}

\textbf{Contexto probable:} sistema de dos poblaciones con respuesta funcional y ejemplos de campo de fase.

\textbf{Transcripcion tentativa:}
\begin{itemize}[leftmargin=*]
    \item Se ve la misma estructura general:
    \[
    \frac{dP}{dt}=rP\left(1-\frac{P}{K}\right)-f(P)D,
    \qquad
    \frac{dD}{dt}=-\gamma D+\delta f(P)D.
    \]
    \item El caso lineal tipo Lotka-Volterra aparece como \(f(P)=P\).
    \item En los dibujos se muestran varios retratos de fase asociados a diferentes respuestas funcionales.
    \item Tambien aparece la idea de que el modelo depende de una respuesta funcional, no solo de un termino bilineal \(PD\).
\end{itemize}

\section*{Foto 19}
\texttt{Fotos/photo\_2026-06-04\_00-13-35.jpg}

\textbf{Contexto probable:} sistemas autonomos en el plano, linealizacion y ejemplos de Jacobiano.

\textbf{Transcripcion tentativa:}
\begin{itemize}[leftmargin=*]
    \item Titulo general: ``Sistemas de ecuaciones diferenciales autonomos''.
    \item Forma del sistema:
    \[
    \frac{dx}{dt}=P(x,y)=a_{11}x+a_{12}y,
    \qquad
    \frac{dy}{dt}=Q(x,y)=a_{21}x+a_{22}y.
    \]
    \item Se escribe
    \[
    F(x,y)=A\binom{x}{y},
    \qquad
    A=\begin{pmatrix}a_{11}&a_{12}\\ a_{21}&a_{22}\end{pmatrix}.
    \]
    \item Desarrollos de Taylor cerca de un punto \(x_0=(x_0,y_0)\) para linearizar alrededor de un equilibrio.
    \item ``Reflejan una caracteristica biologica importante''.
    \item Se dibujan retratos de fase tipo silla, atractor y repulsor.
\end{itemize}

\section*{Foto 20}
\texttt{Fotos/photo\_2026-06-04\_00-13-36.jpg}

\textbf{Contexto probable:} ejemplo de sistema lineal y estabilidad por autovalores reales.

\textbf{Transcripcion tentativa:}
\begin{itemize}[leftmargin=*]
    \item Se trabaja con la linealizacion
    \[
    Y' = DF(x_0)Y = J(x_0)Y.
    \]
    \item En el ejemplo se obtiene un jacobiano diagonal:
    \[
    DF(x_0)=
    \begin{pmatrix}
    2 & 0\\
    0 & 2
    \end{pmatrix},
    \qquad
    DF(x_1)=
    \begin{pmatrix}
    -2 & 0\\
    0 & 2
    \end{pmatrix}.
    \]
    \item Se indica el caso de autovalor doble \(\lambda_{1,2}=2\).
    \item Para \((A-\lambda I)V=0\) se obtiene
    \[
    \det(A-\lambda I)=(2-\lambda)^2.
    \]
    \item Se identifica el caso como una fuente.
\end{itemize}

\section*{Foto 21}
\texttt{Fotos/photo\_2026-06-04\_00-13-37.jpg}

\textbf{Contexto probable:} segundo ejemplo de linealizacion con autovalores distintos y clasificacion de la solucion.

\textbf{Transcripcion tentativa:}
\begin{itemize}[leftmargin=*]
    \item Se considera de nuevo el jacobiano en un equilibrio.
    \item El polinomio caracteristico aparece como
    \[
    p(\lambda)=\det(A-\lambda I)=-(\lambda+2)(\lambda-2),
    \]
    por lo que
    \[
    \lambda_1=2,\qquad \lambda_2=-2.
    \]
    \item Vectores propios:
    \[
    V_1=\binom{1}{0},
    \qquad
    V_2=\binom{0}{1}.
    \]
    \item Soluciones del sistema lineal:
    \[
    y_1(t)=e^{2t}V_1,
    \qquad
    y_2(t)=e^{-2t}V_2.
    \]
    \item La clasificacion indicada en el dibujo es de tipo fuente/silla, segun el equilibrio considerado.
\end{itemize}

\section*{Foto 22}
\texttt{Fotos/photo\_2026-06-04\_00-13-38.jpg}

\textbf{Contexto probable:} referencias bibliograficas y paso del sistema lineal al estudio local de sistemas no lineales.

\textbf{Transcripcion tentativa:}
\begin{itemize}[leftmargin=*]
    \item Se listan referencias:
    \begin{itemize}
        \item Devaney, Hirsch, Smale.
        \item ``Differential equations, dynamical systems and linear algebra''.
        \item ``Dynamical systems''.
    \end{itemize}
    \item Se vuelve a enfatizar:
    \[
    X'=AX \quad \text{en } \mathbb{R}^n,
    \]
    y la solucion general en terminos de exponentiales de matriz.
    \item A la derecha se introduce
    \[
    X' = F(X) \quad \text{en } \mathbb{R}^n,
    \qquad F:\mathbb{R}^n\to\mathbb{R}^n,\ \text{de clase } C^1.
    \]
    \item Problema general: entender la dinamica global de las soluciones de \((*)\).
    \item Dinamica local: dado \(x_0\in\mathbb{R}^n\), estudiar el comportamiento de las soluciones cerca de \(x_0\).
    \item Posibilidades:
    \[
    F(x_0)\neq 0 \quad \text{o} \quad F(x_0)=0.
    \]
\end{itemize}

\section*{Foto 23}
\texttt{Fotos/photo\_2026-06-04\_00-13-39.jpg}

\textbf{Contexto probable:} version ampliada de la dinamica local con ejemplos geometricos.

\textbf{Transcripcion tentativa:}
\begin{itemize}[leftmargin=*]
    \item Se repite la idea de que \(X'=F(X)\) define un campo vectorial en \(\mathbb{R}^n\).
    \item Se muestran trayectorias tipo nodo, espiral y silla.
    \item En el lado derecho se enuncia de nuevo:
    \[
    X_0 \text{ es punto regular si } F(X_0)\neq 0,
    \qquad
    X_0 \text{ es equilibrio si } F(X_0)=0.
    \]
    \item El objetivo es clasificar el comportamiento cerca de esos puntos.
\end{itemize}

\section*{Foto 24}
\texttt{Fotos/photo\_2026-06-04\_00-13-40.jpg}

\textbf{Contexto probable:} continuacion de la dinamica local con retratos de fase y linealizacion.

\textbf{Transcripcion tentativa:}
\begin{itemize}[leftmargin=*]
    \item Se conserva el marco de sistemas no lineales en \(\mathbb{R}^n\).
    \item Se presentan diagramas de trayectorias cercanas a una region regular y a un equilibrio.
    \item La idea de fondo es que el flujo cerca de un punto regular es equivalente a un campo no nulo, mientras que cerca de un equilibrio se requiere linealizacion.
\end{itemize}

\section*{Foto 25}
\texttt{Fotos/photo\_2026-06-04\_00-13-44.jpg}

\textbf{Contexto probable:} linealizacion de un sistema de dimension dos y jacobiano en un equilibrio.

\textbf{Transcripcion tentativa:}
\begin{itemize}[leftmargin=*]
    \item Se escribe
    \[
    X'=F(X), \qquad F \text{ de clase } C^1.
    \]
    \item Para \(X_0=(x_0,y_0)\) equilibrio,
    \[
    J(X_0)=DF(X_0)\in M_{2\times 2}.
    \]
    \item En el caso \(n=2\),
    \[
    J(X)=
    \begin{pmatrix}
    \dfrac{\partial F_1}{\partial x}(x,y) & \dfrac{\partial F_1}{\partial y}(x,y)\\[6pt]
    \dfrac{\partial F_2}{\partial x}(x,y) & \dfrac{\partial F_2}{\partial y}(x,y)
    \end{pmatrix}.
    \]
    \item Evaluando en \(X_0\):
    \[
    J(X_0)=
    \begin{pmatrix}
    \dfrac{\partial F_1}{\partial x}(x_0,y_0) & \dfrac{\partial F_1}{\partial y}(x_0,y_0)\\[6pt]
    \dfrac{\partial F_2}{\partial x}(x_0,y_0) & \dfrac{\partial F_2}{\partial y}(x_0,y_0)
    \end{pmatrix}.
    \]
\end{itemize}

\section*{Foto 26}
\texttt{Fotos/photo\_2026-06-04\_00-13-45.jpg}

\textbf{Contexto probable:} ejemplo linealizado con autovalor doble y clasificacion por vectores propios.

\textbf{Transcripcion tentativa:}
\begin{itemize}[leftmargin=*]
    \item Se trata un ejemplo con
    \[
    x'=x^2-y^2-1,\qquad y'=2y
    \]
    y su jacobiano
    \[
    DF(x)=
    \begin{pmatrix}
    2x & -2y\\
    0 & 2
    \end{pmatrix}.
    \]
    \item Los puntos de equilibrio son
    \[
    X_0=(1,0),\qquad X_1=(-1,0).
    \]
    \item Para el sistema linealizado:
    \[
    Y'=DF(X_0)Y,
    \]
    y en el caso mostrado aparece un autovalor doble \(\lambda=2\).
    \item Se calcula
    \[
    \det(A-\lambda I)=(2-\lambda)^2,
    \]
    con vectores propios coordenados en los ejes.
\end{itemize}

\section*{Foto 27}
\texttt{Fotos/photo\_2026-06-04\_00-13-46.jpg}

\textbf{Contexto probable:} linealizacion de un sistema no lineal y comparacion entre el sistema original y el sistema lineal.

\textbf{Transcripcion tentativa:}
\begin{itemize}[leftmargin=*]
    \item Se reescribe el problema para un equilibrio \(X_0\) y su jacobiano \(J(X_0)\).
    \item Se observa el caso de autovalor doble y la forma
    \[
    (A-\lambda I)V=0.
    \]
    \item La conclusion es que la solucion linealizada tiene el mismo tipo cualitativo de comportamiento cerca del equilibrio.
\end{itemize}

\section*{Foto 28}
\texttt{Fotos/photo\_2026-06-04\_00-13-47.jpg}

\textbf{Contexto probable:} ejemplo con autovalores distintos, fuente, y reconstruccion de soluciones lineales.

\textbf{Transcripcion tentativa:}
\begin{itemize}[leftmargin=*]
    \item Para el jacobiano del equilibrio se obtiene
    \[
    \lambda_1=2,\qquad \lambda_2=-2,
    \]
    con vectores propios
    \[
    V_1=\binom{1}{0},\qquad V_2=\binom{0}{1}.
    \]
    \item Soluciones particulares:
    \[
    y_1(t)=e^{2t}V_1,\qquad y_2(t)=e^{-2t}V_2.
    \]
    \item El retrato de fase mostrado corresponde a una fuente.
\end{itemize}

\section*{Foto 29}
\texttt{Fotos/photo\_2026-06-04\_00-13-48.jpg}

\textbf{Contexto probable:} formalizacion de la solucion general del sistema lineal y base de soluciones por autovalores/vectores propios.

\textbf{Transcripcion tentativa:}
\begin{itemize}[leftmargin=*]
    \item Si \(A\) tiene \(n\) valores propios distintos \(\lambda_1,\dots,\lambda_n\) y vectores propios correspondientes \(v_1,\dots,v_n\), entonces las soluciones
    \[
    X_i(t)=e^{\lambda_i t}v_i
    \]
    son linealmente independientes.
    \item La solucion general es una combinacion lineal de esas soluciones.
    \item Se recuerda que si \(X'=AX\) y \(A v=\lambda v\), entonces
    \[
    X(t)=e^{\lambda t}v
    \]
    es solucion.
\end{itemize}

\section*{Foto 30}
\texttt{Fotos/photo\_2026-06-04\_00-13-50.jpg}

\textbf{Contexto probable:} ejemplo final de linealizacion con sistema no lineal y clasificacion por autovalores.

\textbf{Transcripcion tentativa:}
\begin{itemize}[leftmargin=*]
    \item Se vuelve a mostrar el ejemplo
    \[
    x'=x^2-y^2-1,\qquad y'=2y,
    \]
    con equilibria \((1,0)\) y \((-1,0)\).
    \item Se compara el sistema no lineal con el sistema linealizado
    \[
    Y'=DF(X_0)Y.
    \]
    \item En el lado derecho se ve la descomposicion espectral del sistema lineal y la conclusion de tipo fuente.
\end{itemize}

\section*{Foto 31}
\texttt{Fotos/photo\_2026-06-04\_00-13-51.jpg}

\textbf{Contexto probable:} teorema de Hartman-Grobman y equivalencia topologica local.

\textbf{Transcripcion tentativa:}
\begin{itemize}[leftmargin=*]
    \item ``Teorema de Hartman-Grobman''.
    \item Si \(X_0\) es un punto de equilibrio de
    \[
    X'=F(X),
    \]
    y \(J(X_0)=DF(X_0)\) no tiene valores propios con parte real cero, entonces las soluciones del sistema no lineal cerca de \(X_0\) son topologicamente equivalentes a las del sistema linealizado cerca de \(0\):
    \[
    V'=DF(X_0)V.
    \]
    \item Se introduce un homeomorfismo entre vecindades de \(X_0\) y de \(0\).
\end{itemize}

\section*{Foto 32}
\texttt{Fotos/photo\_2026-06-04\_00-13-56.jpg}

\textbf{Contexto probable:} cierre del teorema de Hartman-Grobman y efectos sobre el retrato de fase.

\textbf{Transcripcion tentativa:}
\begin{itemize}[leftmargin=*]
    \item Se refuerza la idea de que el comportamiento cualitativo cerca de un equilibrio hiperb\'olico se conserva por linealizacion.
    \item Los dibujos acompañan la conclusion con retratos de fase alrededor de nodos, focos y sillas.
\end{itemize}

\section*{Foto 33}
\texttt{Fotos/photo\_2026-06-04\_00-13-57.jpg}

\textbf{Contexto probable:} resumen de dinamica local y clasificacion de puntos regulares y de equilibrio.

\textbf{Transcripcion tentativa:}
\begin{itemize}[leftmargin=*]
    \item Se reitera que el problema general es entender la dinamica global de las soluciones de \(X'=F(X)\).
    \item Dinamica local:
    \[
    \text{dado } X_0\in\mathbb{R}^n,\ \text{como es el comportamiento de las soluciones cerca de }X_0?
    \]
    \item Posibilidades:
    \begin{enumerate}[leftmargin=*]
        \item \(X_0\) es un punto regular, \(F(X_0)\neq 0\).
        \item \(X_0\) es un punto de equilibrio, \(F(X_0)=0\).
    \end{enumerate}
\end{itemize}

\end{document}
